\documentclass[aspectratio=169]{beamer}
\usetheme{Madrid}
\usecolortheme{whale}

\usepackage[utf8]{inputenc}
\usepackage[spanish]{babel}
\usepackage{amsmath,amssymb}
\usepackage{graphicx}
\usepackage{listings}
\usepackage{booktabs}

\lstset{
    language=Python,
    basicstyle=\ttfamily\footnotesize,
    keywordstyle=\color{blue}\bfseries,
    stringstyle=\color{red},
    commentstyle=\color{gray}\itshape,
    breaklines=true,
    showstringspaces=false
}

\title[INT210 - Sesión 0.5]{Sesión 0.5: Listas, Tuplas y Vectores Matemáticos}
\subtitle{Módulo 0.2: Estructuras de Datos Nativas en Python}
\author[Diplomado en Ciencia de Datos]{Cátedra de Inteligencia Artificial y Ciencia de Datos}
\institute[INT210]{Machine Learning From Scratch}
\date{Semana 2}

\begin{document}

\frame{\titlepage}

\begin{frame}{Contenido de la Sesión}
    \tableofcontents
\end{frame}

\section{Introducción y Analogía}

\begin{frame}{La Analogía: El Tren de Carga vs El Mapa en Piedra}
    \begin{columns}
        \column{0.55\textwidth}
        \begin{itemize}
            \item \textbf{Lista (\texttt{list}):} Tren de vagones numerados. Mutable, permite añadir (\texttt{append}), eliminar e intercambiar carga dinámica.
            \item \textbf{Tupla (\texttt{tuple}):} Coordenadas grabadas en piedra. Inmutable, tamaño fijo, acceso instantáneo y seguridad ante modificaciones.
        \end{itemize}
        \column{0.45\textwidth}
        \begin{alertblock}{Principio de Diseño}
            Los pesos y gradientes entrenables se manipulan en listas; las dimensiones fijas y registros de configuración se protegen en tuplas.
        \end{alertblock}
    \end{columns}
\end{frame}

\section{Álgebra Lineal 'From Scratch'}

\begin{frame}{Vectores en $\mathbb{R}^n$ (Joel Grus, Cap. 4)}
    Representamos un vector como \texttt{Vector = List[float]}:
    \begin{columns}
        \column{0.5\textwidth}
        \textbf{Suma Vectorial:}
        \begin{equation*}
            \vec{u} + \vec{v} = [u_1 + v_1, \dots, u_n + v_n]
        \end{equation*}
        \textbf{Multiplicación Escalar:}
        \begin{equation*}
            c \cdot \vec{v} = [c \cdot v_1, \dots, c \cdot v_n]
        \end{equation*}
        \column{0.5\textwidth}
        \textbf{Producto Punto (\textit{Dot}):}
        \begin{equation*}
            \vec{u} \cdot \vec{v} = \sum_{i=1}^n u_i \cdot v_i
        \end{equation*}
        \textbf{Distancia Euclidiana:}
        \begin{equation*}
            d(\vec{u}, \vec{v}) = \sqrt{\sum_{i=1}^n (u_i - v_i)^2}
        \end{equation*}
    \end{columns}
\end{frame}

\section{Aplicaciones en Videojuegos y Arte}

\begin{frame}[fragile]{Colisiones en Motores 3D y Mezcla RGBA}
    \textbf{Detección de Impacto por Distancia al Cuadrado:}
\begin{lstlisting}
def colision_optimizada(pos_a, pos_b, radio_contacto):
    # Evita math.sqrt() para maximizar FPS
    dist_sq = sum((a - b)**2 for a, b in zip(pos_a, pos_b))
    return dist_sq <= radio_contacto**2
\end{lstlisting}
    \vspace{0.2cm}
    \textbf{Mezcla Alfa de Canales RGBA:}
\begin{lstlisting}
r_out = int(a * r_capa + (1.0 - a) * r_fondo)
\end{lstlisting}
\end{frame}

\begin{frame}{Conclusiones}
    \begin{itemize}
        \item Las listas y tuplas proporcionan los bloques constructivos elementales del Álgebra Lineal.
        \item La optimización numérica (como evitar raíces innecesarias) es vital en sistemas de tiempo real.
        \item \textbf{Siguiente Sesión (0.6):} Diccionarios Hash, \texttt{defaultdict} y Análisis de Frecuencia (NLP).
    \end{itemize}
\end{frame}

\end{document}
